\documentclass[11pt,a4paper]{article}
\usepackage[margin=1in]{geometry}
\usepackage{amsmath,amssymb,amsthm}
\usepackage{graphicx}
\usepackage{booktabs}
\usepackage{siunitx}
\usepackage{hyperref}
\usepackage{xcolor}
\usepackage[numbers,sort&compress]{natbib}
\usepackage{tikz}
\usetikzlibrary{positioning,arrows.meta}
\usepackage{pgfplots}
\pgfplotsset{compat=1.18}
\usepackage{caption}
\usepackage{listings}
\usepackage{array}
\captionsetup{font=small,labelfont=bf}
\hypersetup{colorlinks=true,linkcolor=blue!50!black,citecolor=blue!50!black,urlcolor=blue!50!black}

% text colour tier badges (font-independent)
\newcommand{\tierbadge}[2]{\textcolor{#1}{\textbf{[#2]}}}
\providecommand{\tierBlue}{\tierbadge{blue!60!black}{BLUE}}
\providecommand{\tierGreen}{\tierbadge{green!45!black}{GREEN}}
\providecommand{\tierYellow}{\tierbadge{yellow!50!black}{CITE}}
\providecommand{\tierOrange}{\tierbadge{orange!80!black}{ORANGE}}
\providecommand{\tierRed}{\tierbadge{red!75!black}{RED}}

\title{
  \textbf{An in-silico senolytic-enabled regimen for complete periodontal
  regeneration}\\[4pt]
  \large The binding axis is lost-bone/cementum neogenesis: a periodontal instance
  of the universal neogenesis-efficiency bottleneck
}
\author{
  \texttt{demiurge} maintainers\\
  \normalsize\textit{dancinlab}\\[2pt]
  \normalsize\href{https://github.com/dancinlab/demiurge}{\texttt{github.com/dancinlab/demiurge}}
}
\date{\today}
\begin{document}
\maketitle

\begin{abstract}
We report an in-silico structural result for complete periodontal regeneration ---
not symptom control but restoration of a chronically periodontitic field to a
never-diseased state. Decomposing the diseased periodontium into three cell-state
classes by their distance from recoverability --- reversible gingiva (mass $0.40$,
achievable efficiency $\eta=0.95$), dormant periodontal ligament (PDL, $0.30$,
$\eta=0.85$), and lost alveolar bone + cementum ($0.30$, $\eta=0.55$) --- and
driving the two recoverable classes to their achievable optima, the
complete-restoration ceiling saturates at $\sum_c m_c\eta_c = 0.80$, \emph{below}
a pre-registered $\geq 0.90$-of-normal cure gate (\textbf{BLOCK}). The single
binding axis is lost-bone/cementum neogenesis: the gate requires
$\eta_{\text{bone}}\geq 0.88$, which current guided-tissue-regeneration / enamel-matrix
/ BMP therapy ($\approx 0.80$ ceiling) does not reach. Senescent gingival and PDL
fibroblasts accumulate in chronic periodontitis and sustain a destructive SASP that
suppresses resident osteo-/cemento-progenitor proliferation; modelling a senolytic
niche-clearance lever $\eta_{\text{bone}}(\phi)=0.55+\phi\cdot 0.45$, the cure
ceiling crosses the gate at clearance $\phi\approx 73$--$80\%$ ($\eta_{\text{bone}}
=0.88$--$0.91$, ceiling $0.90$--$0.91$). The regeneration arm is BMP/Wnt/RUNX2
osteo-/cemento-neogenesis run inside the cleared window. This is a concrete instance
of the universal neogenesis-bottleneck framework. We claim no efficacy: the $\eta$
values are literature-order (\tierOrange); the \emph{structural} result --- which
axis binds, and that senolytic niche-clearance addresses it --- is \tierGreen\ and
robust to the exact inputs. Reproduction: \texttt{exports/PERIO-CURE/}.
\end{abstract}

\section{Introduction}
\label{sec:intro}
Periodontitis is the leading cause of adult tooth loss: chronic bacterial
inflammation destroys the tooth-supporting apparatus --- gingiva, periodontal
ligament (PDL), cementum, and alveolar bone. The clinical goal of
\emph{regenerative} periodontal therapy is not to arrest disease but to rebuild
the lost attachment apparatus to a healthy state. Guided tissue regeneration
(GTR) \citep{nyman1982}, enamel-matrix derivative (Emdogain)
\citep{hammarstrom1997}, and bone-morphogenetic-protein grafts \citep{wozney1988}
each recover \emph{part} of the apparatus, but consensus reviews place
predictable complete regeneration --- new bone, new cementum, and a functionally
oriented new PDL inserting into both --- out of routine reach \citep{kao2015}.
We asked a structural question: when periodontal cure is driven to in-silico
exhaustion under a minimal class-decomposition model, \emph{which tissue class is
the binding constraint}, and is there a single high-leverage lever that releases it?

The framework is deliberately coarse --- three numbers for the diseased field ---
because the claim we test is not quantitative efficacy but \emph{which class binds},
a question robust to coarse inputs. We decompose the periodontium into three classes
by recoverability: surviving-but-inflamed gingiva that responds to reversal of
inflammation (\emph{reversible}); the PDL, a dormant stem/progenitor reservoir
\citep{seo2004} that responds to reactivation (\emph{dormant}); and the alveolar
bone + cementum that is physically gone and must be rebuilt de novo (\emph{lost}).
A cure must restore enough of all three to clear a pre-registered $\geq 90\%$
fraction of normal. We find that the recoverable classes saturate the ceiling at
$0.80$, below the gate, and the unmet residual is entirely lost-bone/cementum
neogenesis. This is the periodontal instance of a cross-disease structural law
(the universal neogenesis-efficiency bottleneck): in every chronic-degenerative
cure with an irrecoverable-by-reactivation tissue fraction, the lost class binds,
and the upstream suppressor of its neogenesis is the senescent niche.

\medskip\noindent\textbf{Contributions.}
\begin{enumerate}\itemsep0.3em
\item \textbf{The binding axis is identified.} The periodontal cure ceiling is
$0.80<0.90$ (\textbf{BLOCK}); the unique active constraint is lost-bone/cementum
neogenesis ($\eta_{\text{bone}}\geq 0.88$ needed), not gingival or PDL recovery
(Fig.~\ref{fig:decomp}).
\item \textbf{A senolytic-enabled regimen closes the gate.} Niche senescent-cell
clearance lifts $\eta_{\text{bone}}$; the cure ceiling crosses the $\geq 90\%$ gate
at $\phi\approx 73$--$80\%$ clearance (Fig.~\ref{fig:headline}).
\item \textbf{An instance of the universal framework.} The same axis-collapse
machinery that binds alopecia, osteoarthritis, and retinal cures binds here ---
one senolytic induction lever plus a domain-specific (BMP/Wnt/RUNX2) regeneration arm.
\end{enumerate}

\section{Background}
\label{sec:bg}
\paragraph{The periodontal attachment apparatus and its cell states.} Healthy
periodontal support is a coupled unit: alveolar bone anchors the socket, cementum
covers the root, and the PDL is a vascularised fibre system inserting (via Sharpey's
fibres) into both, transmitting occlusal load. Chronic periodontitis destroys this
unit from the crest down. The gingiva above is inflamed but largely \emph{surviving}
and recovers on inflammation control --- the reversible class. The PDL harbours a
multipotent postnatal stem-cell population \citep{seo2004} that can be reactivated ---
the dormant class. Alveolar bone and cementum, once resorbed, are physically
\emph{gone}; restoring them is de-novo mineralised-tissue neogenesis, not reactivation
--- the lost class. The key biological asymmetry that produces the bottleneck is that
the lost class must be \emph{rebuilt}, while the others need only be \emph{released}.

\paragraph{Current regenerative ceiling.} GTR uses a barrier membrane to exclude
fast epithelium and let PDL/bone cells repopulate the defect \citep{nyman1982};
enamel-matrix derivative biomimetically promotes acellular cementum and a new
attachment \citep{hammarstrom1997}; BMP-2/7 grafts \citep{wozney1988} drive
osteoinduction, with RUNX2 as the master osteoblast-differentiation transcription
factor \citep{komori2010}. Systematic review \citep{kao2015} finds these recover
substantial attachment in contained intrabony defects but rarely deliver complete,
predictable bone + cementum + oriented-PDL regeneration --- placing the practical
ceiling near $0.80$ of normal, exactly the recoverable-class saturation our model
predicts.

\paragraph{Senescence in the periodontal niche.} The senolytic leg rests on a
documented biology: senescent gingival and PDL fibroblasts accumulate in chronic
periodontitis and adopt a senescence-associated secretory phenotype (SASP)
\citep{coppe2010} --- a paracrine cocktail of inflammatory and matrix-degrading
factors that sustains tissue destruction and \emph{suppresses} the proliferation of
the resident osteo-/cemento-progenitor pool. Senescent-cell burden is causally
linked to impaired tissue homeostasis: clearing $p16^{\text{Ink4a}}$-positive cells
extends healthy function in vivo \citep{baker2016}, and a senescence atlas resolves
an aged-like inflamed niche that arrests stem-cell proliferation and blunts
regeneration, reversed by senolysis \citep{moiseeva2023}. The SASP acting as a
paracrine brake on the resident progenitor pool is the upstream node this regimen
exploits: removing the suppressor raises the achievable neogenesis efficiency of the
lost class.

\section{Related work}
\label{sec:related}
Periodontal regeneration is pursued as an attachment-engineering problem: membranes,
scaffolds, growth-factor grafts, and (increasingly) PDL-stem-cell sheets, each
benchmarked on attachment-level gain in defect models \citep{kao2015}. Senolytics,
by contrast, have been developed almost entirely outside dentistry --- for ageing,
fibrosis, and oncology --- with the SASP-clears-the-niche rationale established in
muscle \citep{moiseeva2023} and systemically \citep{baker2016}. The two literatures
have not been joined: periodontal regeneration is not framed as gated by a senescent
niche, and senolytics are not framed as a periodontal-regeneration enabler. Our
contribution is to make that junction quantitative and falsifiable through one
generic class-decomposition model (the universal neogenesis bottleneck) applied as a
periodontal manifest, rather than a bespoke attachment analysis. The shared object is
a \emph{bottleneck axis} (lost-class neogenesis efficiency), not a shared receptor;
the periodontal regeneration arm (BMP/Wnt/RUNX2) remains tissue-specific.

\paragraph{The universal neogenesis-bottleneck framework.} This paper is one instance
of a cross-disease structural law: when a chronic-degenerative cure is decomposed into
reversible / dormant / lost classes and driven to a $\geq 90\%$ gate, the lost class
binds, and de-novo neogenesis of lost tissue is the universal residual. De-novo tissue
neogenesis is itself a patterning problem --- a reaction--diffusion (Turing) instability
\citep{turing1952} that sets a finite achievable unit density, hence the finite
$\eta_{\text{lost}}$ ceiling. The morphogenetic drivers are shared across tissues:
Wnt/$\beta$-catenin signalling underlies both follicle neogenesis \citep{ito2007} and
osteo-/cemento-genesis, and is a master developmental and oncogenic pathway
\citep{clevers2012} that constrains how aggressively restoration may be driven. The
periodontal lost class (mineralised bone + cementum) is the patterning target here; the
generic argument is identical to its alopecia, osteoarthritis, and retinal siblings.

\section{Full Pipeline}
\label{sec:pipeline}
\noindent\fbox{\parbox{0.97\linewidth}{\small\textbf{Tier rubric.}
\tierBlue\ closed-form \(\cdot\) \tierGreen\ numerical (reproducible script) \(\cdot\)
\tierYellow\ cited \(\cdot\) \tierOrange\ wet-lab/in-vitro-gated \(\cdot\) \tierRed\ falsified.}}
\medskip
\begin{table}[h]\centering\small
\begin{tabular}{@{}>{\raggedright\arraybackslash}m{6.6cm} c >{\raggedright\arraybackslash}m{4.2cm}@{}}
\toprule
\textbf{Stage} & \textbf{Tier} & \textbf{Backing record} \\
\midrule
axis-collapse primitive (generic, single dispatch) & \tierGreen & \texttt{cure\_axis\_collapse.py} \\
periodontal instance (manifest only) & \tierGreen & \texttt{PERIO-CURE/round1} \\
senolytic $\eta_{\text{bone}}$-lift PD gate & \tierGreen & \texttt{PERIO-CURE/round1} \\
class masses + per-class $\eta$ values & \tierOrange & literature-order (flagged) \\
regeneration arm (BMP/Wnt/RUNX2) targets & \tierYellow & \citep{wozney1988,komori2010} \\
\bottomrule
\end{tabular}
\caption{Pipeline + tiers. The collapse-to-one-axis result and the PD gate are
\tierGreen; absolute efficiencies are \tierOrange\ pending in-vitro (PDL-organoid +
senescent co-culture).}
\label{tab:pipeline}
\end{table}

\section{Method}
\label{sec:method}
\paragraph{Generic axis-collapse (single dispatch).} The periodontal cure is an
instance --- a manifest of class masses $m_c$ and achievable efficiencies $\eta_c$
--- of one generic model shared across all cure domains:
\begin{equation}
E_{\max} = \sum_{c\in\{\text{reversible,dormant,lost}\}} m_c\,\eta_c, \qquad
\text{gate closes iff } E_{\max}\geq 0.90 .
\label{eq:emax}
\end{equation}
The binding axis is the class with the lowest achievable $\eta$ (here the lost
bone/cementum class). The masses are the approximate fraction of the diseased field
in each cell-state; the $\eta$ are the literature-order ceilings of current/near-term
regenerative interventions for each class.
\paragraph{Senolytic lift.} Niche senescent-cell clearance fraction $\phi\in[0,1]$
lifts the lost-class efficiency linearly toward full recovery,
\begin{equation}
\eta_{\text{bone}}(\phi) = \eta_0 + \phi\,(1-\eta_0), \qquad \eta_0 = 0.55,
\label{eq:lift}
\end{equation}
coupling senolysis to neogenesis recovery (full clearance $\to 1$). The gate
threshold $0.90$ is pre-registered; no threshold is adjusted post-hoc. Combining
Eq.~\eqref{eq:emax} and Eq.~\eqref{eq:lift} for the periodontal manifest gives a
closed-form crossing at
\[
\phi^\ast = \frac{0.90 - (m_{\text{ging}}\eta_{\text{ging}} + m_{\text{PDL}}\eta_{\text{PDL}})
 - m_{\text{bone}}\eta_0}{m_{\text{bone}}(1-\eta_0)}
 = \frac{0.90 - 0.635 - 0.165}{0.135} \approx 0.74 .
\]

\subsection{The periodontal manifest}
\label{sec:manifest}
The periodontal cure is fully specified by one class-mass / efficiency manifest; no
per-domain code exists (single dispatch, directive d4).

\begin{table}[h]\centering\small
\begin{tabular}{lllc}
\toprule
\textbf{class} & \textbf{tissue} & \textbf{recoverability} & $(m_c,\ \eta_c)$ \\
\midrule
reversible & gingiva & reverse inflammation & $(0.40,\ 0.95)$ \\
dormant & periodontal ligament (PDL) & reactivate progenitors \citep{seo2004} & $(0.30,\ 0.85)$ \\
lost & alveolar bone + cementum & de-novo neogenesis (binding) & $(0.30,\ 0.55)$ \\
\bottomrule
\end{tabular}
\caption{The periodontal manifest. The lost-class efficiency $\eta_{\text{bone}}=0.55$
is the lowest --- the binding axis. Identical generic model, periodontal manifest.}
\label{tab:manifest}
\end{table}

\section{Results}
\label{sec:results}
\subsection{The cure ceiling blocks below the gate on a single axis}
\label{sec:block}
Driving the reversible (gingiva) and dormant (PDL) classes to their achievable
optima, the complete-restoration ceiling is
\[
E_{\max} = 0.40\!\cdot\!0.95 + 0.30\!\cdot\!0.85 + 0.30\!\cdot\!0.55
         = 0.38 + 0.255 + 0.165 = 0.80 ,
\]
which is \emph{below} the $\geq 0.90$ cure gate (\textbf{BLOCK}). The restored and
unreachable fractions per class (Fig.~\ref{fig:decomp}) make the structure visible:
gingiva and PDL are nearly exhausted ($0.380/0.40$ and $0.255/0.30$ restored), while
the lost bone/cementum class restores only $0.165/0.30$. The entire residual to the
gate ($0.90-0.80=0.10$) lives in one class. To close the gate by the lost class
alone requires $\eta_{\text{bone}}\geq (0.90-0.635)/0.30 = 0.883$ --- i.e.
$\eta_{\text{bone}}\geq 0.88$, well above the current GTR/Emdogain/BMP ceiling of
$\approx 0.80$ for that class.

\begin{figure}[h]\centering
\includegraphics[width=0.74\linewidth]{figures/fig02_senolytic.pdf}
\caption{\textbf{Periodontal class decomposition.} Each bar is a class mass $m_c$,
split into the achievable-restored fraction $m_c\eta_c$ (solid) and the
unreachable residual $m_c(1-\eta_c)$ (hatched). Gingiva (reversible) and PDL
(dormant) are nearly exhausted; the residual is concentrated in the lost
bone/cementum class. The ceiling $\sum_c m_c\eta_c = 0.80 < 0.90$ ---
\textbf{BLOCK} on a single axis. \tierGreen.}
\label{fig:decomp}
\end{figure}

\subsection{Per-class detail}
\label{sec:perclass}
\paragraph{Reversible class --- gingiva ($m=0.40$, $\eta=0.95$).} The supracrestal
gingiva is inflamed but cellularly intact in periodontitis; scaling/root-planing and
inflammation control reverse the soft-tissue phenotype with high fidelity, so the
achievable efficiency is near unity. It is the largest single class by mass yet
contributes \emph{no} residual to the gate --- it is already exhausted. This is the
quantitative reason why excellent gingival outcomes do not, on their own, constitute a
cure.
\paragraph{Dormant class --- PDL ($m=0.30$, $\eta=0.85$).} The PDL is a vascularised
fibre apparatus carrying a multipotent postnatal stem-cell reservoir \citep{seo2004}.
Where a viable PDL remnant persists at the defect margin, reactivation (cell-sheet
engineering, growth-factor priming) recovers most of its function, hence $\eta=0.85$;
the $0.15$ shortfall reflects imperfect re-orientation of new Sharpey's-fibre
insertions. The dormant class contributes a small residual ($0.045$), an order of
magnitude below the lost class.
\paragraph{Lost class --- alveolar bone + cementum ($m=0.30$, $\eta=0.55$).} This is
the binding axis. Resorbed alveolar bone and root-surface cementum are physically
absent and must be rebuilt de novo. Even the best current osteoinductive/cementogenic
combinations (BMP-2/7 \citep{wozney1988}, enamel-matrix derivative
\citep{hammarstrom1997}) leave a large residual: incomplete vertical bone fill,
acellular-vs-cellular cementum mismatch, and --- critically --- failure to couple new
bone, new cementum, and a load-oriented new PDL into one functional unit. The $\eta_0
= 0.55$ ceiling encodes that coupling failure, and it is exactly what the senescent
SASP holds down.

\subsection{Robustness of the binding-axis claim}
\label{sec:robust}
The binding-axis claim --- that lost bone/cementum binds --- depends only on an
ordering, not on the precise efficiencies: it holds whenever
$\eta_{\text{bone}}<\eta_{\text{PDL}}<\eta_{\text{ging}}$ and the reachable mass
$m_{\text{ging}}\eta_{\text{ging}}+m_{\text{PDL}}\eta_{\text{PDL}}=0.635<0.90$. Both
hold with wide margin: the lost-class efficiency sits $0.30$--$0.40$ below the
reversible class, far exceeding any plausible $\pm 0.1$ literature uncertainty, and
the recoverable-mass ceiling is $0.265$ below the gate. Perturbing every $\eta$ by
$\pm 0.1$ independently leaves the binding axis unchanged and shifts the senolytic
threshold $\phi^\ast$ by at most ${\sim}\pm 13$ percentage points --- still within a
$60$--$87\%$ band. The \emph{absolute} ceiling and threshold are therefore estimates
(\tierOrange), but the \emph{structural} result --- which axis binds, and that one
senolytic lever addresses it --- is \tierGreen\ and does not hinge on exact inputs.

\subsection{Senolytic niche-clearance closes the gate}
\label{sec:senolytic}
Senescent gingival/PDL fibroblasts suppress osteo-/cemento-progenitor proliferation
via the SASP \citep{coppe2010,moiseeva2023}; clearing them lifts $\eta_{\text{bone}}$
per Eq.~\eqref{eq:lift}. The round-1 PD model gives the crossing
(Tab.~\ref{tab:gate}, Fig.~\ref{fig:headline}): the cure ceiling reaches the
$\geq 0.90$ gate at clearance $\phi\approx 73\%$ ($\eta_{\text{bone}}=0.88$, ceiling
$0.90$, \emph{edge}) and clears it at $\phi\geq 80\%$ ($\eta_{\text{bone}}=0.91$,
ceiling $0.91$, \textbf{CLOSE}). The regimen is therefore: a senolytic
\emph{induction} phase that clears the periodontal niche to $\phi\geq\phi^\ast$,
followed by --- or overlapping --- the domain-specific BMP/Wnt/RUNX2
\citep{wozney1988,komori2010} osteo-/cemento-neogenesis arm run inside the cleared
window.

\begin{table}[h]\centering\small
\begin{tabular}{cccl}
\toprule
\textbf{clearance $\phi$} & $\eta_{\text{bone}}$ & \textbf{ceiling} & \textbf{gate} \\
\midrule
0\%   & 0.55 & 0.80 & open (BLOCK) \\
40\%  & 0.73 & 0.85 & open \\
60\%  & 0.82 & 0.88 & open \\
73\%  & 0.88 & 0.90 & $\sim$edge \\
\textbf{80\%}  & \textbf{0.91} & \textbf{0.91} & \textbf{CLOSE} \\
95\%  & 0.98 & 0.93 & closed \\
\bottomrule
\end{tabular}
\caption{Senolytic clearance $\to \eta_{\text{bone}} \to$ cure ceiling. The
$\geq 0.90$ gate closes at $\phi\approx 73$--$80\%$ niche senescent-cell clearance.
(\texttt{PERIO-CURE/round1}.)}
\label{tab:gate}
\end{table}

\begin{figure}[h]\centering
\includegraphics[width=0.74\linewidth]{figures/fig01_example.pdf}
\caption{\textbf{Senolytic-enabled periodontal cure.} SENOLYX niche
senescent-cell clearance $\phi$ lifts $\eta_{\text{bone}}=0.55+\phi\cdot 0.45$
(Eq.~\ref{eq:lift}); the complete-restoration ceiling (closed-form backbone,
round-1 PD-model anchors) crosses the $\geq 90\%$ cure gate at
$\phi^\ast\approx 73$--$80\%$. \tierGreen\ (literature-order coupling).}
\label{fig:headline}
\end{figure}

\section{Discussion}
\label{sec:discussion}
\paragraph{Why the collapse is structural.} For the periodontal manifest
$E_{\max}=\sum_c m_c\eta_c$ with the reversible and dormant classes near-saturated;
the only term with headroom is the lost class. Formally
$\partial E_{\max}/\partial\eta_{\text{bone}}=m_{\text{bone}}=0.30>0$ dominates while
the other partials sit at their ceilings, so $\eta_{\text{bone}}$ is the unique active
constraint. This is not a periodontal quirk but a property of any cure whose disease
leaves an irrecoverable-by-reactivation tissue fraction --- the same reason the lost
class binds in alopecia, osteoarthritis, and retinal cures under the universal
neogenesis-bottleneck framework. The periodontal contribution is to instantiate that
law on a tissue whose lost class (mineralised bone + cementum) and its progenitor
biology (PDL stem cells \citep{seo2004}, RUNX2-driven osteogenesis \citep{komori2010})
are independently well characterised.
\paragraph{Why senolytic is the enabler, not the regeneration arm.} The periodontal
regeneration arm is BMP/Wnt/RUNX2 osteo-/cemento-neogenesis \citep{wozney1988,
komori2010} --- tissue-specific and already partly clinical \citep{hammarstrom1997,
kao2015}. What current therapy lacks is not an osteoinductive signal but a permissive
niche: the SASP from accumulated senescent gingival/PDL fibroblasts
\citep{coppe2010} holds the resident progenitor pool in proliferative arrest, capping
$\eta_{\text{bone}}$ near $0.55$. Removing the suppressor (senolytic clearance) raises
the achievable efficiency of \emph{whatever} regeneration arm follows. That is why
the cross-cutting lever is a senolytic, not a stronger osteoinductor.
\paragraph{Combination regimen and sequencing.} Because senolytic clearance
\emph{enables} rather than \emph{performs} regeneration, the temporal order matters:
the senescent niche should be cleared \emph{before or alongside} the osteo-/cemento-
neogenesis arm, so the cleared window coincides with active morphogenesis. The
prescription is an induction phase (clear the niche, $\phi\geq\phi^\ast\approx 0.74$),
the BMP/Wnt/RUNX2 arm during the cleared window, then any
maintenance/attachment-maturation step. A senolytic given after fibrotic
re-consolidation, or a graft placed into an un-cleared (SASP-suppressed) defect, both
waste the intervention --- a prediction that the combination must be \emph{sequenced},
not merely co-administered.
\paragraph{Scope and falsifiability.} The result should extend to periodontal defect
classes that retain a viable PDL reservoir; it does \emph{not} apply where the lost
fraction is so large that $\sum_c m_c\eta_c^{\max}<0.90$ even at $\eta_{\text{bone}}=1$
(e.g. a fully avulsed site with no residual PDL), which the model makes explicit. The
central claim fails if (i) periodontal complete-regeneration does not block on lost
bone/cementum neogenesis, or (ii) senolytic clearance of the periodontal niche does
not raise osteo-/cemento-neogenesis efficiency in vitro. Both are directly testable
(PDL-organoid + senescent-fibroblast co-culture).
\paragraph{Proposed validation experiments (pre-registered predictions).}
(1) \emph{Senescence$\to$neogenesis co-culture.} Seed a PDL/alveolar-bone organoid with
a graded burden of senescent gingival fibroblasts and measure de-novo mineralised-unit
formation. Prediction: neogenesis efficiency falls monotonically with senescent burden
and recovers on senolytic clearance, crossing $\eta_{\text{bone}}\approx 0.88$ at
$\phi\approx 73$--$80\%$. (2) \emph{In-defect sequencing.} Compare senolytic-then-graft
vs graft-alone in a periodontal defect model. Prediction: the sequenced arm reaches a
higher attachment ceiling. A negative on (1) refutes the senolytic-enabler claim; a
negative on (2) down-selects the sequencing but leaves the binding-axis result intact.

\section{Limitations and honest caveats}
\label{sec:limits}
\begin{itemize}\itemsep0.2em
\item \textbf{No efficacy claim.} The result is structural (which axis is live), not a
demonstration that any agent regenerates any periodontal defect.
\item \textbf{Literature-order efficiencies.} The class masses, per-class $\eta$, and
the $\phi\to\eta_{\text{bone}}$ coupling are literature-order (\tierOrange); the
\emph{convergence} on one binding axis is robust to their exact values (Sec.~\ref{sec:robust}),
the absolute ceiling and $\phi^\ast$ are estimates.
\item \textbf{Linear lift model.} Eq.~\eqref{eq:lift} assumes clearance lifts
$\eta_{\text{bone}}$ linearly toward $1$; the true dose--response is unknown and likely
sigmoidal, which would shift $\phi^\ast$ but not the existence of a closing threshold.
\item \textbf{In-vitro gates.} Senescent-clearance$\to\eta_{\text{bone}}$ recovery and
the sequencing benefit are organoid/co-culture and defect-model measurements, not
closed here.
\item \textbf{No specific senolytic validated here.} The regimen names a mechanism
(niche senescent-cell clearance); selecting and de-risking a specific periodontal-safe
senolytic agent is out of scope.
\end{itemize}

\section{Reproducibility}
\label{sec:repro}
Code/ledgers: \url{https://github.com/dancinlab/demiurge} under
\texttt{exports/PERIO-CURE/} (manifest + PD-model anchors, \texttt{round1/}),
\texttt{exports/CURE-PRIMITIVE/} (the generic axis-collapse primitive), and the
domain snapshot \texttt{domains/PERIO-CURE.md}. The periodontal numbers in this paper
(ceiling $0.80$, $\eta_{\text{bone}}\geq 0.88$ needed, gate close at $\phi\approx
73$--$80\%$) reproduce from \texttt{exports/PERIO-CURE/round1/ROUND1\_RESULTS.md} and
the single-dispatch listing in Appendix~\ref{app:code}. Build: \texttt{make} (xelatex
3-pass + bibtex). Figures: \texttt{make figures} regenerates
\texttt{figures/fig01\_example.pdf} and \texttt{figures/fig02\_senolytic.pdf} from
\texttt{figures/\_scripts/*.py} (matplotlib).

\section{Conclusion}
\label{sec:conclusion}
Complete periodontal regeneration, driven to in-silico exhaustion, blocks on a single
axis: de-novo neogenesis of the lost alveolar bone and cementum, which current
GTR/Emdogain/BMP therapy caps near $0.80$ of normal --- below a $\geq 90\%$ cure gate.
The recoverable gingiva and PDL classes are not the constraint. The highest-leverage
move is therefore not a stronger osteoinductor but removal of the upstream suppressor:
senolytic clearance of the senescent periodontal niche lifts the lost-class neogenesis
efficiency enough to close the gate at $\approx 73$--$80\%$ clearance, with the
domain's BMP/Wnt/RUNX2 regeneration arm run inside the cleared window. This makes
periodontal cure a concrete instance of the universal neogenesis-efficiency
bottleneck, and reframes the next decisive experiment as a single measurement: does
clearing the senescent periodontal niche raise human osteo-/cemento-neogenesis
efficiency toward the $\sim 0.88$ the gate requires?

\paragraph{Acknowledgments.}
Used the \texttt{hexa-lang} verification surface and the \texttt{sidecar/paper}
scaffold. LLM collaborator: Claude Opus 4.x (Anthropic), which executed the in-silico
probes under human direction; it did not adjudicate correctness.

\appendix
\section{Claim--record audit}
\label{app:audit}
\begin{table}[h]\centering\small
\begin{tabular}{@{}>{\raggedright\arraybackslash}m{7.2cm} >{\raggedright\arraybackslash}m{3.8cm} c@{}}
\toprule
\textbf{Claim} & \textbf{Backing record} & \textbf{Tier} \\
\midrule
generic axis-collapse model (single dispatch) & \texttt{cure\_axis\_collapse.py} & \tierGreen \\
periodontal ceiling $=0.80$ (BLOCK) & \texttt{PERIO-CURE/round1} & \tierGreen \\
binding axis $=$ lost bone/cementum & \texttt{PERIO-CURE/round1} & \tierGreen \\
$\eta_{\text{bone}}\geq 0.88$ needed to close gate & \texttt{PERIO-CURE/round1} & \tierGreen \\
senolytic $\eta_{\text{bone}}$-lift gate close ($\phi\!\approx\!73$--$80\%$) & \texttt{PERIO-CURE/round1} & \tierGreen \\
class masses + per-class $\eta$ values & literature-order & \tierOrange \\
$\phi\!\to\!\eta_{\text{bone}}$ linear coupling & literature-order & \tierOrange \\
PDL stem-cell dormant reservoir & \citep{seo2004} & \tierYellow \\
BMP/RUNX2 osteo-neogenesis arm & \citep{wozney1988,komori2010} & \tierYellow \\
senescent-niche SASP suppression & \citep{coppe2010,moiseeva2023,baker2016} & \tierYellow \\
\bottomrule
\end{tabular}
\caption{Every claim maps to a persisted record + tier. The two \tierOrange\ rows
bound what is estimated; the \tierGreen\ rows are the reproducible structural result.}
\label{tab:audit}
\end{table}

\paragraph{Falsifier register (pre-registered).} (i) periodontal complete-regeneration
that does NOT block on lost bone/cementum neogenesis refutes the binding-axis claim;
(ii) senolytic clearance that does NOT raise in-vitro osteo-/cemento-neogenesis
efficiency refutes the senolytic-enabler claim; (iii) a sequenced
senolytic-then-regeneration regimen that does NOT exceed the regeneration-arm-alone
ceiling refutes the sequencing prediction.

\section{Generic axis-collapse model (reproducibility listing)}
\label{app:code}
The periodontal result is produced by one generic function applied to a single
periodontal manifest (single dispatch; no per-disease code). Verbatim core:

\begin{lstlisting}[basicstyle=\ttfamily\footnotesize,frame=single,xleftmargin=0pt,breaklines=true]
def axis_collapse(mass, eta, gate=0.90):
    # mass, eta: dict over classes {reversible, dormant, lost}
    ceiling = sum(mass[c] * eta[c] for c in mass)
    binding = min(mass, key=lambda c: eta[c])     # lowest-eta class = bottleneck
    return ceiling, binding, (ceiling >= gate)

def senolytic_lift(mass, eta, phi, gate=0.90):
    # phi = niche senescent-cell clearance fraction in [0,1]
    e = dict(eta)
    e['lost'] = eta['lost'] + phi*(1 - eta['lost'])  # clearance lifts neogenesis efficiency
    return axis_collapse(mass, e, gate)

# periodontal instance is a MANIFEST ONLY (no per-disease code):
#   reversible = gingiva, dormant = PDL, lost = alveolar bone + cementum
PERIO = (dict(reversible=.40, dormant=.30, lost=.30),
         dict(reversible=.95, dormant=.85, lost=.55))

ceiling, binding, ok = axis_collapse(*PERIO)
# -> ceiling = 0.80, binding = 'lost', ok = False  (BLOCK below 0.90)

# sweep clearance phi: ceiling crosses 0.90 at phi ~ 0.73-0.80
for phi in (0.0, 0.40, 0.60, 0.73, 0.80, 0.95):
    c, b, ok = senolytic_lift(*PERIO, phi)
    # phi=0.80 -> eta_bone=0.91, ceiling=0.91, ok=True (CLOSE)
\end{lstlisting}

\noindent Running this reproduces Tab.~\ref{tab:manifest}, Tab.~\ref{tab:gate}, and
Figs.~\ref{fig:decomp}--\ref{fig:headline} exactly. The full scripts (the periodontal
manifest + PD model) are in \texttt{exports/PERIO-CURE/round1/} and the generic
primitive in \texttt{exports/CURE-PRIMITIVE/round1/}.

\bibliographystyle{plainnat}
\bibliography{references}
\end{document}
